\section{Adaptation Cost: What Each Change Requires}
\label{sec:adaptation}

A deployed decision component serves a pool of candidate models, a set of decisions and a set of deployment constraints, such as an offload ratio, a token budget, a latency target or a risk target. Each of these can change after deployment, which gives three axes of change: 1) new candidate models (\axisM), 2) new types of decisions (\axisD) and 3) new constraints (\axisC).

\noindent\textbf{Reusing the model is not free.} Many changes leave the shared model's weights unchanged. They may still require additional data and settings, such as a model profile, question text, a threshold or a small head fitted for one decision. Preparing the data and settings for each change has an \emph{adaptation cost} with four parts: 1) outcome labels, 2) calls to profile a model, 3) unlabeled warmup samples and 4) the specification text a person must write. We count newly acquired resources separately from reused ones. For example, reusing an existing calibration set costs no new labels but still relies on labeled evidence.

\noindent\textbf{A sequence of changes.} Consider a support service that classifies tickets and routes difficult cases to an LLM. Adding a ticket category is a new decision: the team writes its definition and checks it on a small set of labeled tickets, while the shared model stays frozen. Adding an LLM is a new model: the team records its outputs on a reference set and supplies the resulting profile. Lowering the offload ratio is a new constraint: the team picks a new threshold from representative warmup scores. Tightening the permitted error rate is also a new constraint, but it needs a risk-control procedure. The team may reuse its existing labeled set if the traffic, the question text and the model version have not changed since the set was labeled.

\begin{table*}[!t]
\centering
\caption{What a change on each axis requires, by family. \axisM: a new candidate model; \axisD: a new decision; \axisC: a new deployment constraint, either an offload ratio or a risk target.}
\label{tab:requirements}
\footnotesize
\setlength{\tabcolsep}{4pt}
\renewcommand{\arraystretch}{1.05}
\begin{tabularx}{\textwidth}{@{}>{\raggedright\arraybackslash\hspace{0pt}}p{3.4cm}>{\raggedright\arraybackslash\hspace{0pt}}X>{\raggedright\arraybackslash\hspace{0pt}}X>{\raggedright\arraybackslash\hspace{0pt}}X@{}}
\toprule
\tblhead \thd{Family (examples)} & \thd{\axisM{} New model} & \thd{\axisD{} New decision} & \thd{\axisC{} New constraint} \\
\midrule
\tblgroup{4}{Routers and gates}
Profile routers (UniRoute~\cite{jitkrittum2025universal}; kNN~\cite{li2025rethinking}) & A profile built from the model's graded outputs on a reference set & Not supported: one fixed routing question & A cost weight applied to predicted per-model errors; labeled evidence for a risk target \\
Route-as-text routers (Arch-Router~\cite{tran2025arch}) & An edit to the human-maintained policy-to-model table & A new route description; routes only & Not addressed \\
Agreement pivots (Yue et~al.~\cite{yue2024large}; CARGO~\cite{chen2026routing}) & A capable cheap model with comparable answers, and a retuned agreement rule & One fixed accept-or-escalate question & An agreement threshold or probabilistic offloading pivot (CARGO calibrates the latter on an unlabeled warmup batch per offload ratio); labeled evidence for a risk target \\
\midrule
\tblgroup{4}{Decisions specified in the input}
Zero-shot classifiers (NLI~\cite{yin2019benchmarking}; GLiClass~\cite{stepanov2025gliclass}) & Not used for model selection & New label strings & A threshold on label scores; labeled evidence for a risk target \\
Policy-conditioned guards (Llama Guard~\cite{inan2023llama}; DynaGuard~\cite{hoover2025dynaguard}) & Not used for model selection & New policy text; Boolean verdicts, plus the violated categories in Llama Guard & A threshold on the probability of an unsafe or failing verdict; labeled evidence for a risk target \\
Rubric judges (Prometheus~2~\cite{kim2024prometheus2}; G-Eval~\cite{liu2023geval}) & Not a router; can grade a new model's outputs to build its profile & A new rubric; scores and pairwise preferences & A threshold on scores; labeled evidence for a risk target \\
\midrule
\tblgroup{4}{Typed decision models}
Jev and open implementations (Table~\ref{tab:implementations}) & Candidate-specific evidence, such as a profile built from graded outputs, supplied in the state or used by the surrounding router & New question text, and labeled examples to test it & Warmup scores for an offload ratio; a validated acceptance rule with labeled evidence for a risk target \\
\bottomrule
\end{tabularx}
\end{table*}


\noindent\textbf{What each family needs.} Table~\ref{tab:requirements} lists what a change on each axis requires, family by family. Typed decision models accept a new decision as question text, without changing the shared model. Testing the new decision still takes labeled examples. The other two axes need further information. In the routing approaches reviewed here, a new candidate model needs candidate-specific evidence, such as a profile (Section~\ref{sec:routing}). An offload ratio needs representative unlabeled scores. A risk target needs a validated acceptance rule and labeled evidence.

\noindent\textbf{Guarantees need labeled evidence.} The methods that set or certify an acceptance rule use labeled data. Temperature scaling fits one parameter on labeled data to reduce overconfidence~\cite{guo2017calibration}, but it carries no finite-sample guarantee. Selective prediction chooses a confidence threshold on held-out labels so that the accepted decisions meet a risk level with high probability~\cite{geifman2017selective}. Conformal prediction turns a labeled calibration set into prediction sets with a coverage guarantee under exchangeability~\cite{angelopoulos2021gentle}. Conformal risk control extends the guarantee to a bounded monotone loss, such as a miss rate~\cite{angelopoulos2024conformal}. Trust or Escalate applies a labeled threshold of this kind to a cascade of LLM judges, with a guarantee of human agreement~\cite{jung2025trust}. These guarantees hold on the distribution they were calibrated on. They need not hold within every subgroup or after the model or the population changes, and calibration is known to degrade under dataset shift~\cite{ovadia2019can}. A changed question, population or model version can therefore require new labeled evidence.
